\section*{Chapter \ref{ch:g7:moon-phases} --- Phases of the Moon}

\begin{solution}{exo:g7:moon-phases:1}
Exactly half, always --- the sunward half. A \omterm{def:g7:moon-phases:phases}{phase} is only the
view: the fraction of that lit half which happens to face the
Earth.
\end{solution}

\begin{solution}{exo:g7:moon-phases:2}
New, \omterm{def:g7:moon-phases:phases}{waxing} \omterm{ex:g2:moon-in-the-sky:shapes}{crescent}, first quarter, \omterm{def:g7:moon-phases:phases}{waxing} gibbous, full,
\omterm{def:g7:moon-phases:phases}{waning} gibbous, last quarter, \omterm{def:g7:moon-phases:phases}{waning} \omterm{ex:g2:moon-in-the-sky:shapes}{crescent} --- about $29.5$
days around. \omterm{def:g7:moon-phases:phases}{Waxing} is the growing half, from new to full.
\end{solution}

\begin{solution}{exo:g7:moon-phases:3}
New moon: you face the lamp with the ball toward it --- its lit
half turned away from you. \omterm{ex:g2:moon-in-the-sky:shapes}{Full moon}: lamp at your back, ball
before you --- its lit half turned fully toward you.
\end{solution}

\begin{solution}{exo:g7:moon-phases:4}
Any two: the Earth's \omterm{def:g1:light-and-shadows:shadow}{shadow} points away from the Sun and could
only reach a \omterm{def:g2:moon-in-the-sky:moon}{Moon} at full position; a \omterm{ex:g2:moon-in-the-sky:shapes}{crescent} \omterm{def:g2:moon-in-the-sky:moon}{Moon} stands
sunward of us, nowhere near that \omterm{def:g1:light-and-shadows:shadow}{shadow}; and a round \omterm{def:g1:light-and-shadows:shadow}{shadow}'s
edge could never carve the gibbous double curve.
\end{solution}

\begin{solution}{exo:g7:moon-phases:5}
The new moon is dark to us because its \emph{own} night side
faces us --- no \omterm{def:g1:light-and-shadows:shadow}{shadow} of ours involved, and the far side is in
full sunshine. The \omterm{def:g7:shadows-eclipses:lunar}{lunar eclipse} is the full \omterm{def:g2:moon-in-the-sky:moon}{Moon} crossing the
Earth's real \omterm{prop:g7:shadows-eclipses:umbra}{umbra}: our \omterm{def:g1:light-and-shadows:shadow}{shadow} genuinely falling on it, briefly.
\end{solution}

\begin{solution}{exo:g7:moon-phases:6}
The full \omterm{def:g2:moon-in-the-sky:moon}{Moon} stands opposite the Sun. As the spinning Earth
carries the Sun down below the western horizon, the opposite
point of the sky --- where the full \omterm{def:g2:moon-in-the-sky:moon}{Moon} waits --- must be
rising in the east: \omterm{def:g1:day-and-night:daynight}{sunset} and full-moonrise are one event seen
both ways.
\end{solution}

\begin{solution}{exo:g7:moon-phases:7}
The new moon travels with the Sun: at midnight both stand under
your feet, on the far side of the Earth. A new moon at midnight
is as impossible as the Sun at midnight.
\end{solution}

\begin{solution}{exo:g7:moon-phases:8}
Away from the \omterm{def:g1:day-and-night:daynight}{sunset} point --- the \omterm{ex:g2:moon-in-the-sky:shapes}{crescent}'s bow bends toward
the Sun, horns away. The lit side must face the Sun, and the
horns are the lit side's retreating edges: geometry forbids them
to aim sunward.
\end{solution}

\begin{solution}{exo:g7:moon-phases:9}
Between first quarter and full --- about three-eighths of the
way around, on the evening side. It rose mid-afternoon. Five
nights later: \omterm{ex:g2:moon-in-the-sky:shapes}{full moon} --- rising at \omterm{def:g1:day-and-night:daynight}{sunset}, up all night.
\end{solution}

\begin{solution}{exo:g7:moon-phases:10}
\omterm{def:g7:shadows-eclipses:solar}{Solar eclipse}: only at new moon --- Sun--Moon--Earth in line.
Lunar: only at \omterm{ex:g2:moon-in-the-sky:shapes}{full moon} --- Sun--Earth--Moon in line. Rare even
then, because the \omterm{def:g2:moon-in-the-sky:moon}{Moon}'s tilted \omterm{def:g6:sun-earth-moon:axisorbit}{orbit} usually carries it above
or below the exact line: \omterm{def:g7:moon-phases:phases}{phase} is necessary, alignment is not
guaranteed.
\end{solution}

\begin{solution}{exo:g7:moon-phases:11}
The Earth shows \omterm{def:g7:moon-phases:phases}{phases} opposite to ours, on the same clock:
when we see new moon, the lunar astronauts see a \emph{full
Earth} --- our whole daylit side facing them; at our \omterm{ex:g2:moon-in-the-sky:shapes}{full moon}
they see a new Earth. One geometry, read from the other chair.
\end{solution}

\begin{solution}{exo:g7:moon-phases:12}
Path: sunlight strikes the Earth, the bright Earth diffuses it
moonward, the \omterm{def:g2:moon-in-the-sky:moon}{Moon}'s night side diffuses a little back to our
eyes --- sunlight twice forwarded. At our \omterm{ex:g2:moon-in-the-sky:shapes}{crescent} \omterm{def:g7:moon-phases:phases}{phases} the
\omterm{def:g2:moon-in-the-sky:moon}{Moon} sees a nearly \emph{full Earth} (the \omterm{def:g7:moon-phases:phases}{opposite-phase} rule):
the night side is floodlit by the brightest Earth possible ---
strongest earthshine. Toward our \omterm{ex:g2:moon-in-the-sky:shapes}{full moon} the \omterm{def:g2:moon-in-the-sky:moon}{Moon} sees a thin
\omterm{ex:g2:moon-in-the-sky:shapes}{crescent} Earth: the floodlight dwindles, and \omterm{def:g7:moon-phases:phases}{full-phase} glare
drowns what remains.
\end{solution}

\begin{solution}{pb:g7:moon-phases:1}
\textbf{1.} First quarter near night $7$; \omterm{ex:g2:moon-in-the-sky:shapes}{full moon} near night
$15$; last quarter near night $22$ (the cycle closing near
night $29$--$30$).
\textbf{2.} \omterm{ex:g2:moon-in-the-sky:shapes}{Full moon}: rises at \omterm{def:g1:day-and-night:daynight}{sunset}. First quarter: rises
near midday. Last quarter: rises near midnight. New moon: rises
and sets with the Sun.
\textbf{3.} The young \omterm{ex:g2:moon-in-the-sky:shapes}{crescent} stands only a little behind the
Sun on the \omterm{def:g6:sun-earth-moon:axisorbit}{orbit}: it hangs in the west after \omterm{def:g1:day-and-night:daynight}{sunset} and soon
follows the Sun down. The \omterm{ex:g2:moon-in-the-sky:shapes}{full moon} stands opposite: as the Sun
sets in the west, opposite means rising in the east.
\textbf{4.} Any \omterm{def:g7:moon-phases:phases}{phase} sharing part of the Sun's shift: the
\omterm{def:g7:moon-phases:phases}{waxing} gibbous through afternoon, and above all the \omterm{def:g7:moon-phases:phases}{waning}
\omterm{def:g7:moon-phases:phases}{phases} --- last quarter and \omterm{def:g7:moon-phases:phases}{waning} \omterm{ex:g2:moon-in-the-sky:shapes}{crescent} ride high in
morning skies.
\textbf{5.} \omterm{def:g7:moon-phases:phases}{Waxing} gibbous: the ball three-quarters lit means
you stand between first quarter and full --- lamp behind your
shoulder, ball ahead and to the side.
\textbf{6.} They are seeing the ball's \emph{own} night side ---
the half turned from the lamp. The Earth-you casts a \omterm{def:g1:light-and-shadows:shadow}{shadow} too,
but it stretches behind you, away from the lamp; the ball is
nowhere in it except at the full-moon position --- and then only
if exactly aligned: an eclipse, not a \omterm{def:g7:moon-phases:phases}{phase}.
\textbf{7.} The \omterm{def:g2:moon-in-the-sky:moon}{Moon} turns exactly once per \omterm{def:g6:sun-earth-moon:axisorbit}{orbit}, so the same
half faces the Earth forever --- rabbit, face and all; the far
side waited for a camera to fly around.
\textbf{8.} Spin and \omterm{def:g6:sun-earth-moon:axisorbit}{orbit} keep perfect step: as the \omterm{def:g2:moon-in-the-sky:moon}{Moon}
advances along its path, it turns just enough to keep the same
face pointed at us --- like a dancer circling a partner while
always facing them: truly spinning (the room sees every side),
yet showing the partner one face only.
\textbf{9.} Rising at midnight is the last-quarter schedule ---
the \omterm{def:g7:moon-phases:phases}{waning} side: last quarter (the lit half faces the morning
Sun's direction).
\textbf{10.} Errors: a \omterm{ex:g2:moon-in-the-sky:shapes}{crescent} stands near the Sun, so it can
never ride high at midnight; horns point \emph{away} from the
Sun, never at it; and \omterm{def:g1:day-and-night:daynight}{sunset} colors at midnight are their own
impossibility. Three strikes.
\textbf{11.} No: a \omterm{def:g7:shadows-eclipses:solar}{solar eclipse} requires the \omterm{def:g2:moon-in-the-sky:moon}{Moon} between Sun
and Earth --- the new-moon position. Tonight's \omterm{def:g2:moon-in-the-sky:moon}{Moon} stands at
full, on the opposite side: the two events cannot share a
night.
\textbf{12.} For example: ``One fact runs the month: the \omterm{def:g2:moon-in-the-sky:moon}{Moon}
is always half lit, and we circle-watch that half from
changing sides. Buried this month: the Earth's \omterm{def:g1:light-and-shadows:shadow}{shadow} has
nothing to do with \omterm{def:g7:moon-phases:phases}{phases} --- it visits only at rare eclipses.
And take the clock home: every \omterm{def:g7:moon-phases:phases}{phase} keeps its own rising hour
--- full at \omterm{def:g1:day-and-night:daynight}{sunset}, quarters at midday and midnight --- so the
\omterm{def:g2:moon-in-the-sky:moon}{Moon}'s shape tells you its schedule.''
\end{solution}
